\chapter{Raspodijeljeni sustavi i tolerancija na greške}
\label{ch:raspodijeljeni}

Poziv udaljenom sustavu može završiti na više načina nego lokalni poziv.
Odgovor može kasniti ili izostati, iako je udaljeni sustav obradio zahtjev.
Zbog toga sučelje mora odrediti što učiniti kada ishod nije poznat i što se
događa ako se isti zahtjev pošalje ponovno. Ti se problemi pojavljuju i u
fiskalizacijskom tijeku iz poglavlja \ref{ch:studija}.

% ------------------------------------------------------------
\section{Raspodijeljeni sustav i djelomični otkaz}
\label{sec:djelomicni-otkaz}

Waldo i suradnici razlikuju lokalni i udaljeni poziv
\cite{waldo1994note}. Lokalni poziv ostaje unutar jednog adresnog prostora.
Udaljeni poziv prelazi u drugi adresni prostor, često na drugom računalu.
Pozivatelj pritom poznaje sučelje udaljene komponente, ali ne upravlja njezinim
izvođenjem.

Autori navode četiri razlike između lokalnog i udaljenog poziva. Udaljeni
poziv traje dulje, nema izravan pristup istoj memoriji, može djelomično
otkazati i odvija se istodobno s radom drugih komponenti. Posebno su važni
djelomični otkaz i istodobnost jer ih nije moguće sakriti samom tehničkom
izvedbom poziva.

Kod lokalnog poziva pozivatelj i pozvana funkcija rade u istom procesu. Ako
proces stane, staju oboje. Kod udaljenog poziva jedna strana može prestati
raditi dok druga nastavlja. Primjerice, primatelj može obraditi zahtjev nakon
što je pošiljatelj prestao čekati odgovor. Nijedna strana tada nema potpunu
sliku događaja. Waldo i suradnici to ograničenje opisuju ovako:

\begin{quote}
„U raspodijeljenom sustavu otkaz mrežne veze ne razlikuje se od otkaza
procesora s druge strane te veze.“\footnote{„\emph{In a distributed
system, the failure of a network link is indistinguishable from the
failure of a processor on the other side of that link.}“
\cite{waldo1994note}} \end{quote}

Izostanak odgovora ne otkriva gdje je nastao problem. Zahtjev možda nije
stigao do primatelja. Primatelj je možda prestao raditi prije obrade ili je
zahtjev obradio, ali nije poslao odgovor. Moguće je i da se izgubio samo
odgovor. Sučelje zato mora propisati postupanje kada ishod zahtjeva nije
poznat.

Autori problem prikazuju na redu čekanja sa zadanim sučeljem. Klijent nakon
\strani{isteka vremena}{timeout} može uvesti
\strani{ponavljanje}{retry}. Isti se zahtjev tada može obraditi više puta.
Primatelj može prepoznati ponavljanje samo ako poruka sadrži identifikator
koji ostaje isti u svakom pokušaju. Poruke iz primjera nemaju takav
identifikator, a zadano sučelje ne dopušta njegovo dodavanje. Pouzdaniji red
čekanja ne može nadoknaditi podatak koji nedostaje u poruci. Zato pri
oblikovanju svake operacije treba odrediti učinak njezina ponavljanja:

\begin{quote}
„Pri oblikovanju svake operacije mora se postaviti pitanje: što se događa
ako klijent odluči ponoviti tu operaciju s potpuno istim parametrima kao
prije?“\footnote{„\emph{In the design of any operation, the question has
to be asked: what happens if the client chooses to repeat this operation
with the exact same parameters as previously?}“ \cite{waldo1994note}}
\end{quote}

Isto pitanje vrijedi za fiskalizacijsko sučelje. Njegova su pravila zadana
propisom pa ih pojedini sudionik ne može sam promijeniti.

% ------------------------------------------------------------
\section{Pojmovi: kvar, pogreška i otkaz}
\label{sec:fault-error-failure}

Avizienis i suradnici razlikuju kvar, pogrešku i otkaz
\cite{avizienis2004taxonomy}. Kvar (\emph{fault}) je uzrok problema, poput
propusta u programu. Taj kvar može stvoriti pogrešno unutarnje stanje, što
autori nazivaju pogreškom (\emph{error}). Otkaz (\emph{failure}) nastaje tek
kada usluga zbog te pogreške više ne radi kako je propisano. Kvar je aktivan
dok proizvodi pogrešku, a inače je uspavan.

Primjer to jasno razdvaja. Propust u kodu može pogrešno izračunati ukupan
iznos računa. Propust je kvar, a pogrešan iznos u memoriji pogreška. Ako
sustav pošalje takav račun, nastao je otkaz usluge. Ako prije slanja otkrije
i ispravi iznos, pogreška nije postala vidljiva korisniku i nema otkaza.
Ovaj rad ponajprije analizira takva vidljiva odstupanja u radu usluge.

Problem ne mora biti samo u programu. Program može točno slijediti
specifikaciju, a ipak dati neprihvatljiv rezultat. Tada je kvar u
specifikaciji. Primjeri su izostavljeno pravilo, netočna pretpostavka i
nedosljedne upute. Poglavlje \ref{ch:studija} zato provjerava i pravila koja
određuju fiskalizacijski tijek, a ne samo moguću programsku izvedbu.

Kvar može biti trajan ili prolazan. Trajni kvar ostaje prisutan, primjerice
zbog neispravne konfiguracije. Prolazni kvar traje ograničeno, primjerice
zbog kratkog prekida veze. Kvar je čvrst ako se može ponovno izazvati pod
istim uvjetima. Ako se ne može pouzdano ponoviti, naziva se neuhvatljivim.

Ta razlika određuje ima li smisla pokušati ponovno. Novi pokušaj može uspjeti
nakon prolaznog kvara, ali neće ukloniti trajno neispravnu konfiguraciju.
Odjeljak \ref{sec:failure-modes} provjerava može li pošiljatelj iz odgovora
fiskalizacijskog sučelja prepoznati tu razliku.

% ------------------------------------------------------------
\section{Klasifikacija otkaza}
\label{sec:klasifikacija-otkaza}

Avizienis i suradnici razvrstavaju otkaze prema tome što je pošlo po zlu i
može li korisnik to uočiti \cite{avizienis2004taxonomy}. Te su dvije podjele
najvažnije za ovaj rad.

Sadržajni otkaz znači da je usluga vratila pogrešan podatak. Vremenski otkaz
znači da je ispravan podatak stigao prerano ili prekasno. Ako usluga prestane
slati poruke, riječ je o zaustavljanju. Tihi otkaz poseban je slučaj u kojem
korisnik ne primi ni rezultat ni obavijest o problemu.

Otkaz je signaliziran ako sustav korisniku javi da nešto nije uspjelo. Ako
obavijest izostane, otkaz je nesignaliziran. Sustav pritom može prijaviti
problem koji ne postoji ili propustiti postojeći problem. Sudionik ne može
pokrenuti odgovarajući oporavak ako ne zna što se dogodilo.

Ponovni pokušaj ima smisla samo ako isti kvar vjerojatno neće nastupiti opet.
Gray pretpostavlja da se program nakon vraćanja u početno stanje često može
oporaviti ponavljanjem operacije \cite{gray1986whystop}:

\begin{quote}
„Pretpostavljam da postoji sličan fenomen u softveru: većina proizvodnih
programskih kvarova je meka. Ako se stanje programa vrati na početno i
operacija ponovi, operacija najčešće neće pasti drugi
put.“\footnote{„\emph{I conjecture that there is a similar phenomenon in
software -- most production software faults are soft. If the program state
is reinitialized and the failed operation retried, the operation will
usually not fail the second time.}“ \cite{gray1986whystop}}
\end{quote}

Gray čvrste kvarove naziva \emph{Bohrbug}, a meke \emph{Heisenbug}. Meke je kvarove teško
ponoviti jer samo ispitivanje može promijeniti uvjete u kojima nastaju. Gray
svoju pretpostavku podupire analizom dnevnika grešaka, ali ne određuje koliko
često ponavljanje uspijeva. Ponovni pokušaj zato nije jamstvo uspjeha.

% ------------------------------------------------------------
\section{Istek vremena i otkrivanje otkaza}
\label{sec:detekcija-otkaza}

Chandra i Toueg opisuju detektor otkaza kao dio sustava koji prati druge
procese \cite{chandra1996failuredetectors}. Svaki proces vodi vlastiti popis
procesa za koje sumnja da su otkazali. Sumnja može biti pogrešna i poslije se
povući. Zato dva procesa u istom trenutku ne moraju imati isti popis.

Potpunost znači da će detektor na kraju označiti svaki proces koji je
otkazao. Točnost znači da neće pogrešno označavati ispravne procese. Jaka
točnost zahtijeva da nijedan ispravan proces nikada ne bude označen kao
otkazao. U stvarnom sustavu to je teško postići jer se spor proces izvana
može ponašati jednako kao proces koji je prestao raditi.

Istek vremena zato nije dokaz da je primatelj otkazao. On samo znači da je
pozivatelj prestao čekati. Zahtjev možda nije stigao, možda se još obrađuje,
a možda je već obrađen. Sustav mora razlikovati takav nepoznat ishod od
odgovora koji izričito potvrđuje neuspjeh.

% ------------------------------------------------------------
\section{Jamstva dostave i potvrda ishoda}
\label{sec:semantika-dostave}

Jamstvo dostave govori koliko puta sustav može obraditi istu poruku. Kod
dostave najviše jednom poruka se obrađuje jednom ili se izgubi. Kod dostave
barem jednom sustav ponavlja slanje dok ne dobije potvrdu, pa se poruka može
obraditi više puta. Dostava točno jednom zahtijeva da se poruka obradi jedanput
i samo jedanput.

Ako pošiljatelj ne ponavlja slanje, poruka se neće obraditi više puta, ali se
može izgubiti. Ako šalje sve dok ne primi potvrdu, smanjuje mogućnost gubitka,
ali otvara mogućnost ponovljene obrade. Pritom treba razlikovati tri događaja:

\begin{itemize}[nosep, leftmargin=1.4cm]
  \item poruka je stigla do primatelja,
  \item primatelj je proveo traženu radnju,
  \item oba sudionika znaju ishod radnje.
\end{itemize}

Paradoks dvaju generala pokazuje granicu takvog dogovora
\cite{gray1978notes}. Dva generala moraju potvrditi da će napasti istodobno,
ali se svaki glasnik može izgubiti. Svaka potvrda zato traži novu potvrdu da
je stigla. Ako se posljednja poruka izgubi, jedan general ne zna je li dogovor
potvrđen. Dodavanje još jedne poruke samo premješta isti problem na kraj
protokola.

Taj se rezultat odnosi na protokol s unaprijed ograničenim brojem poruka.
Gray zatim opisuje dvofazni protokol koji može slati neograničen broj poruka.
Takav protokol može postići dogovor, ali nije zajamčeno da će uvijek završiti.
Zaseban formalni rezultat pokazuje da u asinkronom sustavu nije moguće
zajamčiti konsenzus ako može otkazati barem jedan proces
\cite{fischer1985impossibility}.

Isto se može dogoditi pri slanju računa. Primatelj obradi zahtjev i pošalje
potvrdu, ali se potvrda izgubi. Primatelj zna ishod, a pošiljatelj ga ne zna.
Paradoks dvaju generala ovdje služi samo kao objašnjenje te granice. Iz njega
ne slijedi da je dostava točno jednom uvijek nemoguća ni da se svaki zahtjev
mora obraditi više puta.

U praksi se dostava barem jednom često spaja s obradom koja podnosi
ponavljanje. Helland opisuje primatelja koji pamti obrađene poruke. Kada istu
poruku primi ponovno, ne ponavlja radnju nego vraća spremljeni odgovor
\cite{helland2007lifebeyond}.

Operacija je idempotentna ako njezino ponavljanje ne mijenja konačni učinak.
Naredba „postavi status na plaćeno“ ima to svojstvo. Naredba „uvećaj naplaćeni
iznos za 100 eura“ nema ga jer svako izvođenje ponovno uvećava iznos. Ipak,
isti učinak nije dovoljan da pošiljatelj sazna što se dogodilo. Primatelj mora
vratiti raniji odgovor ili omogućiti zaseban dohvat statusa.

Saltzer, Reed i Clark objašnjavaju da mreža ne može sama osigurati funkciju
koja ovisi o podacima poznatima samo aplikaciji \cite{saltzer1984endtoend}.
To vrijedi i za prepoznavanje ponovljenih zahtjeva:

\begin{quote}
„Čak i ako mreža potiskuje duplikate, sama aplikacija može nehotice
proizvesti duplicirane zahtjeve u vlastitim postupcima ponavljanja nakon
otkaza. Komunikacijskom sustavu ta udvostručenja na razini aplikacije
izgledaju kao različite poruke, pa ih on ne može
potisnuti.“\footnote{„\emph{even if the network suppresses duplicates, the
application itself may accidentally originate duplicate requests, in its
own failure/retry procedures. These application level duplications look
like different messages to the communication system, so it cannot suppress
them}“ \cite{saltzer1984endtoend}} \end{quote}

Mreža može potvrditi da je prenijela poruku, ali ne i da je odredišna
aplikacija provela traženu radnju. Tu potvrdu može dati samo aplikacija koja
je zahtjev obradila. Poglavlje \ref{ch:studija} provjerava daje li
fiskalizacijsko sučelje dovoljno podataka za utvrđivanje ishoda.

Gilbert i Lynch opisuju drugo ograničenje, poznato kao CAP
\cite{gilbert2002brewer}. Kratica dolazi od engleskih naziva za usklađenost,
dostupnost i otpornost na razdvajanje mreže: \emph{consistency},
\emph{availability} i \emph{partition tolerance}. Kada mreža razdvoji
čvorove repliciranog sustava, sustav ne može za svaki zahtjev istodobno
jamčiti odgovor i strogo usklađene podatke. Mora privremeno odbiti dio
zahtjeva ili dopustiti da čvorovi daju različite rezultate. I kada nema
razdvajanja, manjem kašnjenju može se dati prednost pred trenutačnom
usklađenošću \cite{abadi2012consistency}.

CAP se ne primjenjuje izravno na ovaj rad jer fiskalizacijski tijek nije
replicirana baza podataka. Ovdje je bitan jednostavniji problem: isti dokument
obrađuju dva odvojena kanala, a ne postoji zajednički protokol koji bi obje
obrade dovršio kao jednu cjelinu.
